\documentclass[../../../include/open-logic-section]{subfiles}

\begin{document}

\olfileid{sth}{ordinals}{iso}
\olsection{هم‌ریختی‌های ترتیبی}

برای توضیح اینکه خوش‌ترتیبی \emph{تا چه اندازه} استوار است، با معرفی روشی
برای مقایسهٔ خوش‌ترتیب‌ها آغاز می‌کنیم.

\begin{defn}
یک \emph{خوش‌ترتیب} زوج~$\tuple{A,<}$ است به‌طوری‌که~$<$، $A$ را
خوش‌ترتیب کند. خوش‌ترتیب‌های~$\tuple{A,<}$ و~$\tuple{B,\lessdot}$
\emph{هم‌ریختِ ترتیبی‌اند} اگر و تنها اگر یک~!!a{bijection}
$f \colon A \to B$ وجود داشته باشد به‌طوری‌که: $x<y$ اگر و تنها اگر
$f(x) \lessdot f(y)$. در این حالت می‌نویسیم
$\ordeq{\tuple{A,<}}{\tuple{B,\lessdot}}$ و می‌گوییم~$f$ یک
\emph{هم‌ریختی ترتیبی} است.
\end{defn}
\noindent
در ادامه، برای اختصار، به‌جای «هم‌ریختی‌های ترتیبی» صرفاً از
«هم‌ریختی‌ها» سخن خواهیم گفت. از دید شهودی، هم‌ریختی‌ها
!!{bijection}های حافظ ساختارند. در ادامه چند واقعیت ساده دربارهٔ
هم‌ریختی‌ها می‌آید.

\begin{lem}\ollabel{isoscompose}
ترکیب هم‌ریختی‌ها هم‌ریختی است؛ یعنی اگر
$f \colon A \to B$ و~$g \colon B \to C$ هم‌ریختی باشند، آنگاه
$(g \circ f) \colon A \to C$ هم‌ریختی است.
\end{lem}
\begin{prob}
	\olref[sth][ordinals][iso]{isoscompose} را اثبات کنید.
\end{prob}
\begin{proof}
به‌عنوان تمرین واگذار می‌شود.
\end{proof}
\begin{cor}\ollabel{ordisoisequiv}
	$\ordeq{X}{Y}$ یک رابطهٔ هم‌ارزی است.
\end{cor}
\begin{prop}\ollabel{ordisounique}
اگر~$\tuple{A,<}$ و~$\tuple{B,\lessdot}$ خوش‌ترتیب‌هایی هم‌ریخت باشند،
هم‌ریختی میان آن‌ها یکتاست.
\end{prop}

\begin{proof}
فرض کنید~$f$ و~$g$ هم‌ریختی‌هایی~$A \to B$ باشند. نتیجه را با استقرا،
یعنی با استفاده از~\olref[wo]{propwoinduction}، اثبات می‌کنیم.
$a\in A$ را ثابت بگیرید و (به‌عنوان فرض استقرا) فرض کنید
$(\forall b<a)f(b)=g(b)$. همچنین~$x\in B$ را ثابت بگیرید.

اگر~$x \lessdot f(a)$، آنگاه~$f^{-1}(x)<a$ و بنابراین
$g(f^{-1}(x)) \lessdot g(a)$، زیرا~$f$ و~$g$ هم‌ریختی‌اند. اما چون
$f^{-1}(x)<a$، بنا بر فرض ما
$x=f(f^{-1}(x))=g(f^{-1}(x))$. پس~$x \lessdot g(a)$. به همین ترتیب، اگر
$x \lessdot g(a)$، آنگاه~$x \lessdot f(a)$.

با تعمیم، $(\forall x \in B)(x \lessdot f(a) \liff x \lessdot g(a))$.
بنا بر~\olref[sfr][rel][ord]{prop:extensionality-strictlinearorders} از
این نتیجه می‌شود که~$f(a)=g(a)$. پس با
\olref[wo]{propwoinduction}، $(\forall a \in A)f(a)=g(a)$.
\end{proof}\noindent
این تا اندازه‌ای نشان می‌دهد که خوش‌ترتیب‌ها استوارند. اما برای ادامهٔ
توضیح بهتر است نمادگذاری بیشتری معرفی کنیم.

\begin{defn}
اگر~$\tuple{A,<}$ خوش‌ترتیبی با~$a\in A$ باشد، بگذارید
$A_a = \Setabs{x \in A}{x<a}$. می‌گوییم~$A_a$ یک \emph{قطعهٔ آغازین}
سره از~$A$ است (و اجازه می‌دهیم خود~$A$ قطعهٔ آغازین ناسره‌ای از~$A$
باشد). بگذارید~$<_a$ تحدید~$<$ به این قطعهٔ آغازین باشد؛ یعنی
$\funrestrictionto{\mathord{<}}{A_a^2}$.
\end{defn}
\noindent
با این نمادگذاری می‌توانیم بیان و اثبات کنیم که هیچ خوش‌ترتیبی با هیچ‌یک
از قطعه‌های آغازین سرهٔ خود هم‌ریخت نیست.

\begin{lem}\ollabel{wellordnotinitial}
اگر~$\tuple{A,<}$ خوش‌ترتیبی با~$a\in A$ باشد، آنگاه
$\ordneq{\tuple{A,<}}{\tuple{A_a,<_a}}$.
\end{lem}

\begin{proof}
برای برهان خلف، فرض کنید~$f \colon A \to A_a$ یک هم‌ریختی باشد. چون~$f$
تناظری یک‌به‌یک است و~$A_a \subsetneq A$، با استفاده
از~\olref[wo]{wo:strictorder} بگذارید~$b\in A$ یک~!!{element}~$<$-کمینهٔ~$A$ باشد
که~$b\neq f(b)$. نشان خواهیم داد
$(\forall x \in A)(x<b \liff x<f(b))$؛ آنگاه بنا
بر~\olref[sfr][rel][ord]{prop:extensionality-strictlinearorders} نتیجه
می‌شود~$b=f(b)$ و برهان خلف کامل می‌شود.

فرض کنید~$x<b$. بنا بر انتخاب~$b$، داریم~$x=f(x)$. و چون~$f$ هم‌ریختی
است، $f(x)<f(b)$. پس~$x<f(b)$.

فرض کنید~$x<f(b)$. چون~$f$ هم‌ریختی است، $f^{-1}(x)<b$ و بنا بر انتخاب
$b$، $f^{-1}(x)=x$. پس~$x<b$.
\end{proof}

نتیجهٔ بعدی ما، به‌بیان تقریبی، نشان می‌دهد یک «قطعهٔ آغازین» از هم‌ریختی
نیز هم‌ریختی است:

\begin{lem}\ollabel{wellordinitialsegment}
فرض کنید~$\tuple{A,<}$ و~$\tuple{B,\lessdot}$ خوش‌ترتیب باشند. اگر
$f \colon A \to B$ یک هم‌ریختی و~$a\in A$ باشد، آنگاه
$\funrestrictionto{f}{A_a} : A_a \to B_{f(a)}$ یک هم‌ریختی است.
\end{lem}

\begin{proof}
چون~$f$ هم‌ریختی است:
	%چون $f$ هم‌ریختی است، $b<a$ اگر و تنها اگر $f(b)\lessdot f(a)$، پس
\begin{align*}
	\funimage{f}{A_a} &= \funimage{f}{\Setabs{x \in A}{x<a}}\\
	&= \Setabs{y \in B}{f^{-1}(y)<a} \\
	&= \Setabs{y \in B}{y \lessdot f(a)} \\
	&=B_{f(a)}
\end{align*}
و~$\funrestrictionto{f}{A_a}$ ترتیب را حفظ می‌کند، زیرا~$f$ آن را حفظ
می‌کند.
\end{proof}

دو نتیجهٔ بعدی نشان می‌دهند خوش‌ترتیب‌ها همواره قابل مقایسه‌اند:

\begin{lem}\ollabel{lemordsegments}
فرض کنید~$\tuple{A,<}$ و~$\tuple{B,\lessdot}$ خوش‌ترتیب باشند. اگر
$\ordeq{\tuple{A_{a_1},<_{a_1}}}{\tuple{B_{b_1},\lessdot_{b_1}}}$ و
$\ordeq{\tuple{A_{{a_2}},<_{a_2}}}{\tuple{B_{{b_2}},
\lessdot_{b_2}}}$، آنگاه
${a_1}<{a_2} \text{ اگر و تنها اگر }{b_1}\lessdot{b_2}$.
\end{lem}

\begin{proof}
جهت \emph{چپ به راست} را اثبات می‌کنیم؛ جهت دیگر مشابه است. فرض کنید هم
$\ordeq{\tuple{A_{a_1},<_{a_1}}}{\tuple{B_{b_1},
\lessdot_{b_1}}}$ و هم~$\ordeq{\tuple{A_{{a_2}},
<_{a_2}}}{\tuple{B_{{b_2}},\lessdot_{b_2}}}$، و
$f \colon A_{{a_2}} \to B_{{b_2}}$ هم‌ریختی ما باشد. فرض کنید
${a_1}<{a_2}$؛ آنگاه بنا بر~\olref{wellordinitialsegment}،
$\ordeq{\tuple{A_{a_1},<_{a_1}}}{\tuple{B_{f({a_1})},
\lessdot_{f({a_1})}}}$. پس
$\ordeq{\tuple{B_{b_1},\lessdot_{b_1}}}{\tuple{B_{f({a_1})},
\lessdot_{f({a_1})}}}$، و بنابراین بنا بر~\olref{wellordnotinitial}،
${b_1}=f({a_1})$. اکنون چون برد~$f$ برابر~$B_{{b_2}}$ است،
${b_1}\lessdot{b_2}$.
\end{proof}

\begin{thm}\ollabel{thm:woalwayscomparable}
برای هر دو خوش‌ترتیب، یکی با قطعه‌ای آغازین (که لزوماً سره نیست) از دیگری
هم‌ریخت است.
\end{thm}

\begin{proof}
فرض کنید~$\tuple{A,<}$ و~$\tuple{B,\lessdot}$ خوش‌ترتیب باشند. با استفاده
از جداسازی، بگذارید
\[
	f = \Setabs{\tuple{a,b} \in A \times B}{
		\ordeq{\tuple{A_a,<_a}}{\tuple{B_b,\lessdot_b}}}.
\]
بنا بر~\olref{lemordsegments}، برای هر
$\tuple{a_1,b_1},\tuple{a_2,b_2}\in f$، $a_1<a_2$ اگر و تنها اگر
$b_1\lessdot b_2$. پس~$f \colon \dom{f} \to \ran{f}$ یک هم‌ریختی است.

اگر~$a_2\in\dom{f}$ و~$a_1<a_2$، آنگاه بنا
بر~\olref{wellordinitialsegment}، $a_1\in\dom{f}$؛ پس~$\dom{f}$ قطعه‌ای
آغازین از~$A$ است. به همین ترتیب، $\ran{f}$ قطعه‌ای آغازین از~$B$ است.
برای برهان خلف، فرض کنید هر دو قطعهٔ آغازین \emph{سره} باشند. آنگاه
بگذارید~$a$ یک~!!{element}~$<$-کمینهٔ~$A\setminus\dom{f}$ باشد، به‌طوری‌که
$\dom{f}=A_a$؛ و بگذارید~$b$ یک~!!{element}~$\lessdot$-کمینهٔ
$B\setminus\ran{f}$ باشد، به‌طوری‌که~$\ran{f}=B_b$. پس
$f \colon A_a \to B_b$ یک هم‌ریختی است، و در نتیجه
$\tuple{a,b}\in f$، که تناقض است.
\end{proof}

\end{document}
